\documentclass[11pt]{article}
\usepackage[margin=1in]{geometry}
\usepackage{amsmath,amssymb,amsthm,booktabs,enumitem,hyperref,natbib}
\hypersetup{
 colorlinks=true,
 linkcolor=blue,
 citecolor=blue,
 urlcolor=blue,
 pdftitle={Maximal divergence of the leximin and leximax core selections in exact games},
 pdfauthor={Maxwell Flanagan}
}

\newtheorem{theorem}{Theorem}
\newtheorem{lemma}[theorem]{Lemma}
\newtheorem{proposition}[theorem]{Proposition}
\newtheorem{corollary}[theorem]{Corollary}
\newtheorem{conjecture}[theorem]{Conjecture}
\newtheorem{observation}[theorem]{Observation}
\newtheorem{definition}[theorem]{Definition}
\newtheorem{remark}[theorem]{Remark}
\newcommand{\core}{\operatorname{Core}}
\newcommand{\Lmin}{\operatorname{Lmin}}
\newcommand{\Lmax}{\operatorname{Lmax}}
\newcommand{\one}{\mathbf 1}
\newcommand{\supp}{\operatorname{supp}}

\title{Maximal divergence of the leximin and leximax\\
core selections in exact games}
\author{Maxwell Flanagan}
\date{Draft of August 6, 2026}

\begin{document}
\maketitle

\begin{abstract}
A singleton-normalized exact transferable-utility game carries two canonical
egalitarian core selections: the leximin allocation $\Lmin$ and the
leximax-min allocation $\Lmax$.  They coincide on every exact game with at
most four players.  We measure their separation by
$\rho(v)=\|\Lmin(v)-\Lmax(v)\|_1/v(N)$, which never exceeds $2$, and show
that $2$ is the exact supremum.  The witnesses form an explicit two-parameter
family $G(p,k)$ of two-generator lower-envelope games with
$\rho(G(p,k))=2(p(k-1)-1)/(k(p+1))$ and the deficit identity
$2-\rho(G(p,k))=2n/(k(p+1))$, $n=p+1+k$; the balanced split gives
$\rho\ge 2-8n/(n^2-(n\bmod 2))$ at every $n\ge6$, so the ratio exceeds $1$
from eight players on.  The mechanism is a single family of tight
$(p+1)$-player caps that both selections saturate: leximin reads a cap as a
ceiling and piles the residual onto one player, leximax-min reads it as a
budget and spreads the residual over half the players, so the two supports
become asymptotically disjoint.  The concentrating selection is the leximin
one.  On the two-generator class we prove a dual description of the core,
$\core(v)=\{x\ge0:x(N)=T,\ x(S)\le\max(g(S),h(S))\}$, deduce a one-coalition
upper bound on the leximin floor, and exhibit the smallest game at which the
mirror bound for the leximax-min ceiling fails.  At $n=5$ the family formula
is not sharp: the true two-generator maximum is $1/2$, proved here by hand
from four core vertices.  We also close one natural route to a matching upper
bound, by a seven-player game in which smoothing a generator toward
uniformity collapses $\rho$ from $1/5$ to $0$.
\end{abstract}

\section{Introduction}

Let $(N,v)$ be a transferable-utility game on $N=\{1,\dots,n\}$ with
\[
 \core(v)=\{x\in\mathbb R^N: x(N)=v(N),\ x(S)\ge v(S)\ \text{for all}\ S\subseteq N\},
\]
where $x(S)=\sum_{i\in S}x_i$.  The game is \emph{exact} if every coalition
inequality is attained by some core allocation, and \emph{singleton-normalized}
if $v(\{i\})=0$ for every $i$; on that domain every core allocation is
nonnegative.  Exact games are the natural home for egalitarian core
selections, and two such selections are canonical.  The \emph{leximin}
allocation $\Lmin(v)$ maximizes the increasing rearrangement of $x$
lexicographically over the core; the \emph{leximax-min} allocation
$\Lmax(v)$ minimizes the decreasing rearrangement lexicographically over the
core.  Both are single-valued on a nonempty compact convex set.  The first
raises the poorest as far as possible, the second lowers the richest as far
as possible.

On small games the two agree.  The companion paper \citep{Flanagan2026}
proves that every singleton-normalized exact game on at most four players is
an exact partition game at $\Lmin$; by the established implication from exact
partition to Lorenz domination \citep{LlerenaMauri2017}, $\Lmin$ then
Lorenz-dominates every core allocation, and a Lorenz-greatest core point is
simultaneously leximin- and leximax-optimal, so
$\Lmin(v)=\Lmax(v)$ for every singleton-normalized exact game with $n\le4$.
Our complete enumerations at $n\le4$ reconfirm this rigidity as a control
(Section~\ref{sec:fixedn}).  The adjacent large-core chain --- exact games on
at most four players have large cores \citep{Biswas1999}, and on large-core
games the procedural egalitarian solution coincides with $\Lmax$
\citep{Dietzenbacher2021,ArinKuipersVermeulen2003} --- pins only the $\Lmax$
side: it does not deliver the coincidence, and indeed the companion
repository exhibits a five-player exact game with a large core on which
$\Lmin\ne\Lmax$.  The companion's lower-envelope realization technique is
used again here.

The question of this paper is what happens after the rigidity breaks.  Write
\[
 \rho(v)=\frac{\|\Lmin(v)-\Lmax(v)\|_1}{v(N)},
 \qquad
 s(n)=\sup\bigl\{\rho(v):\ v\ \text{exact on $n$ players},\ v(N)>0\bigr\},
\]
the supremum running over singleton-normalized games.
Since $\Lmin(v)$ and $\Lmax(v)$ are nonnegative with the same total $v(N)$,
we have $\rho\le2$ always, and $\rho=2$ would mean disjoint supports.  The
rigidity says $s(n)=0$ for $n\le4$.  How large can $s(n)$ get?

\paragraph{Results.}
The answer is: as large as arithmetically possible.

\begin{itemize}[nosep]
\item $\sup_n s(n)=2$, and at fixed $n$, $s(n)\ge 2-8n/(n^2-(n\bmod2))$,
      which is $2-8/n$ for even $n$ (Theorem~\ref{thm:sup}); in particular
      $s(n)>1$ from $n=8$ on, and the bound is attained on the two-generator
      class at $n=6,7,8$ by complete enumeration.
\item The witnesses form a two-parameter family $G(p,k)$ whose two
      allocations are computed in closed form (Theorem~\ref{thm:family}),
      with the deficit identity $2-\rho(G(p,k))=2n/(k(p+1))$
      (Corollary~\ref{cor:deficit}).
\item On the two-generator class the core has a dual description as a system
      of upper caps (Theorem~\ref{thm:duality}), from which the leximin floor
      inherits a one-coalition certificate (Theorem~\ref{thm:floor}); the
      mirror statement for the leximax-min ceiling is false, and we give the
      smallest counterexample (Observation~\ref{obs:Mside}).
\item $n=5$ is anomalous: the family formula gives $1/3$ but the true
      two-generator maximum is $1/2$, proved by hand
      (Theorem~\ref{thm:n5}), and the $n=5$ extremal class is qualitatively
      unlike the $n\ge6$ one.
\end{itemize}

\paragraph{The mechanism, in one sentence.}
In $G(p,k)$ a single family of tight caps forces $p+1$ players to share only
$p$; leximin reads each cap as a ceiling, so the floor sticks at $p/(p+1)$
and the entire residual piles onto one player, while leximax-min reads the
same cap as a budget, loading half the players up to $p(k-1)/k$ and thereby
squeezing the capped players down to $1/k$ --- so the two supports become
asymptotically disjoint and $\rho\to2$.  It is worth stating the direction
explicitly, because it inverts the usual reading: the \emph{concentrating}
selection is the leximin one.  In $G(p,k)$ the leximin allocation gives one
player a share $p(k-1)/((p+1)k)\to1$ of $v(N)$ while everyone else sits just
below $1$; it is nevertheless leximin-optimal, because nothing better is
feasible for the poorest.  Both selections are driven into extreme shapes by
the same constraint, in opposite directions.

\paragraph{What is not proved.}
An upper bound matching Theorem~\ref{thm:sup} at fixed $n$ is open, and
Section~\ref{sec:obstruction} shows that one natural route to it is closed.
Section~\ref{sec:scope} separates the proved statements from the enumerated
and the conjectured ones, and lists the open problems.

\section{Preliminaries}\label{sec:prelim}

Throughout, $N=\{1,\dots,n\}$, all games are singleton-normalized, and
$T:=v(N)>0$.  Vectors are indexed by $N$, and $\supp(x)=\{i:x_i\ne0\}$.

\subsection{Exact games as lower envelopes}

The following generation fact is the only construction used in this paper.
It is standard; see \citet{CsokaHeringsKoczy2011} and
\citet{StudenyKratochvil2022} for the surrounding theory of the exact cone,
and \citet{Flanagan2026} for the same statement in the companion.

\begin{lemma}[lower-envelope generation]\label{lem:envelope}
Let $g_1,\dots,g_r\in\mathbb R_+^N$ have a common total
$g_1(N)=\dots=g_r(N)=T$, and set $v(S)=\min_{j}g_j(S)$.  Then $v$ is exact,
$v(N)=T$, every $g_j$ lies in $\core(v)$, and $v$ is singleton-normalized as
soon as every coordinate vanishes in some $g_j$.  Conversely, every exact
singleton-normalized game is of this form: one may take $g_1,\dots,g_r$ to be
the vertices of $\core(v)$.
\end{lemma}

\begin{proof}
Each $g_j$ satisfies $g_j(S)\ge v(S)$ by definition of the minimum and
$g_j(N)=T=v(N)$, so $g_j\in\core(v)$; and the inequality at $S$ is attained
by whichever $g_j$ realizes the minimum, so $v$ is exact.  The singleton
statement is immediate from $v(\{i\})=\min_j (g_j)_i$.  Conversely, if $v$ is
exact then $v(S)=\min_{x\in\core(v)}x(S)$ for every $S$, and the minimum of a
linear functional over the polytope $\core(v)$ is attained at a vertex; the
vertices are nonnegative because $v$ is singleton-normalized, and all have
total $T$.
\end{proof}

We call a game generated as in Lemma~\ref{lem:envelope} by $r$ allocations an
\emph{$r$-generator} game.  Part of what follows is specific to $r=2$.

\subsection{The two selections}

For $x\in\mathbb R^N_+$ let $x^\uparrow$ be its increasing and
$x^\downarrow$ its decreasing rearrangement.

\begin{definition}
$\Lmin(v)$ is the allocation in $\core(v)$ whose increasing rearrangement is
lexicographically largest; $\Lmax(v)$ is the allocation in $\core(v)$ whose
decreasing rearrangement is lexicographically smallest.
\end{definition}

Both exist and are unique on a nonempty compact convex set: existence is
compactness applied stage by stage, and uniqueness follows by comparing two
putative optima with their midpoint, which is feasible by convexity and
strictly better at the first stage where they differ.  We write
$\mu(v)=\min_i\Lmin(v)_i$ for the \emph{leximin floor} and
$M(v)=\max_i\Lmax(v)_i$ for the \emph{leximax-min ceiling}; equivalently
$\mu=\max_{x\in\core}\min_i x_i$ and $M=\min_{x\in\core}\max_i x_i$.

\subsection{Normalization and the ratio}

Two conventions have to be stated honestly, because neither selection is
translation-equivariant: both are applied to the raw payoff vector $x$, not
to $x-(v(\{i\}))_{i\in N}$, and on a game that is not singleton-normalized
these are different objects --- the leximax-min allocation in particular
changes under the shift.  We therefore fix the domain once and for all: all
games are singleton-normalized, and both selections are computed on the
normalized game.  No claim below asserts invariance of $\Lmin$ or $\Lmax$
under adding singleton values.

Scaling, by contrast, is harmless: $\core(cv)=c\core(v)$ for $c>0$ and both
selections are positively homogeneous, so $\Lmin(cv)=c\Lmin(v)$ and
$\Lmax(cv)=c\Lmax(v)$.  Hence $\rho(v)=\|\Lmin(v)-\Lmax(v)\|_1/v(N)$ is
scale-invariant, which is what makes it comparable across games of different
size.  We record the trivial ceiling and the form of the deficit.

\begin{lemma}[ceiling and overlap]\label{lem:ceiling}
For any singleton-normalized exact $v$ with $T=v(N)>0$,
\[
 \rho(v)=2\Bigl(1-\frac1T\sum_{i\in N}\min\bigl(\Lmin(v)_i,\Lmax(v)_i\bigr)\Bigr)
 \ \le\ 2,
\]
with $\rho(v)=2$ if and only if $\Lmin(v)$ and $\Lmax(v)$ have disjoint
supports.
\end{lemma}

\begin{proof}
For nonnegative $x,y$ with $x(N)=y(N)=T$ one has
$|x_i-y_i|=x_i+y_i-2\min(x_i,y_i)$, and summing gives
$\|x-y\|_1=2T-2\sum_i\min(x_i,y_i)$.  The overlap term is nonnegative and
vanishes exactly when no coordinate is positive in both.
\end{proof}

So $\rho\to2$ is precisely asymptotic support disjointness of the two
egalitarian selections, and the whole content of the paper is that this can
be forced.

\section{The two-generator class}\label{sec:twogen}

Throughout this section $v$ is the lower envelope of exactly two nonnegative
allocations $g,h$ with common total $T$, so that $v(S)=\min(g(S),h(S))$ for
every $S$ (including $S=N$, where both sides equal $T$).

\begin{lemma}[support disjointness]\label{lem:disjoint}
If $v$ is singleton-normalized then $A:=\supp(g)$ and $B:=\supp(h)$ are
disjoint, and every player outside $A\cup B$ is null, receiving $0$ in every
core allocation.  Deleting null players changes neither $\Lmin$, $\Lmax$, nor
$T$, so we may and do assume $A\cup B=N$ and $A\cap B=\emptyset$.
\end{lemma}

\begin{proof}
$v(\{i\})=\min(g_i,h_i)=0$ forces $g_i=0$ or $h_i=0$, which is exactly
$A\cap B=\emptyset$.  If $i\notin A\cup B$ then $g_i=h_i=0$, so
$v(N\setminus\{i\})=\min\bigl(g(N\setminus\{i\}),h(N\setminus\{i\})\bigr)=T$;
any $x\in\core(v)$ has $x(N\setminus\{i\})\ge T$ and $x(N)=T$ with $x\ge0$,
hence $x_i=0$.  Both selections assign a null player $0$, so deletion is
harmless.
\end{proof}

\begin{definition}
For $S\subseteq N$ put $f(S):=\max\bigl(g(S),h(S)\bigr)=\max\bigl(g(S\cap A),h(S\cap B)\bigr)$.
\end{definition}

\subsection{A dual description of the core}

\begin{theorem}[two-generator core duality]\label{thm:duality}
\[
 \core(v)=\bigl\{x\ge0:\ x(N)=T\ \text{and}\ x(S)\le f(S)\ \text{for all }S\subseteq N\bigr\}.
\]
\end{theorem}

\begin{proof}
For $x$ with $x(N)=T$ and any $S$, the inequality $x(S)\ge v(S)$ is
equivalent to $x(N\setminus S)\le T-v(S)$.  It therefore suffices to prove
\[
 T-v(S)=f(N\setminus S)\qquad\text{for every }S\subseteq N.  \tag{1}
\]
Since $g$ and $h$ are modular with $g(N)=h(N)=T$,
\[
 T-v(S)=T-\min\bigl(g(S),h(S)\bigr)
 =\max\bigl(T-g(S),\,T-h(S)\bigr)
 =\max\bigl(g(N\setminus S),\,h(N\setminus S)\bigr)
 =f(N\setminus S).
\]
Replacing $S$ by $N\setminus S$ in (1) turns the family
$\{x(S)\ge v(S)\}$ into the family $\{x(S)\le f(S)\}$.  Finally $x\ge0$ is
exactly the family of singleton inequalities $x(\{i\})\ge v(\{i\})=0$.
\end{proof}

Identity (1) uses only modularity and the common total; the disjointness of
Lemma~\ref{lem:disjoint} is not needed for it, and enters only to normalize
the presentation of $f$.  Note also that $S=N$ gives $f(N)=T$, so the cap at
$N$ is implied by efficiency, and $S=\emptyset$ gives $f(\emptyset)=0$.

In fact the duality is fully general: for \emph{any} singleton-normalized
game, setting $f(S):=T-v(N\setminus S)$, the same complementation involution
gives $\core(v)=\{x\ge0:x(N)=T,\ x(S)\le f(S)\ \forall S\}$.  The
two-generator content of Theorem~\ref{thm:duality} is the computation
$f=\max(g,h)$.  The general form is used in Section~\ref{sec:n5closed}.

\subsection{The leximin floor}

\begin{theorem}[leximin floor: the easy direction]\label{thm:floor}
Let
\[
 \lambda:=\min_{\emptyset\ne S\subseteq N}\frac{f(S)}{|S|}
 =\min_{(A',B')\ne(\emptyset,\emptyset)}
   \frac{\max\bigl(g(A'),h(B')\bigr)}{|A'|+|B'|},
\]
the second minimum running over $A'\subseteq A$, $B'\subseteq B$.  Then
$\mu\le\lambda$.
\end{theorem}

\begin{proof}
$\Lmin(v)$ is a core allocation, so by Theorem~\ref{thm:duality}, for every
nonempty $S$,
$|S|\,\mu\le\Lmin(v)(S)\le f(S)$, i.e.\ $\mu\le f(S)/|S|$.  Take the minimum
over $S$.
\end{proof}

The bound is a \emph{one-coalition certificate}: a single tightest cap
already pins the floor from above.  Empirically it is also exact.

\begin{conjecture}[leximin floor formula]\label{conj:floor}
$\mu=\lambda$ for every singleton-normalized two-generator exact game.
\end{conjecture}

\paragraph{Evidence and reduction.}
Conjecture~\ref{conj:floor} holds in $586$ of $586$ random two-generator
games at $n=4,5,6,7$ (sample sizes $120$, $218$, $165$, $83$), with all
arithmetic exact.  What remains to prove is a clean linear-programming
statement.  Equality means that the system
\[
 y\ge0,\qquad y(N)=T-n\lambda,\qquad y(S)\le f'(S):=f(S)-\lambda|S|
\]
is feasible; indeed, given such $y$, the allocation $x=y+\lambda\one$
satisfies $x(N)=T$, $x\ge\lambda\one\ge0$ and $x(S)\le f(S)$, hence lies in
the core with floor at least $\lambda$.  Note $f'\ge0$ by definition of
$\lambda$.  By linear-programming duality,
\[
 \max\{y(N):\ y\ge0,\ y(S)\le f'(S)\ \forall S\}
 =\min\Bigl\{\textstyle\sum_S w_Sf'(S):\ w\ge0,\ \sum_{S\ni i}w_S\ge1\ \forall i\Bigr\},
\]
the minimum running over fractional covers.  Taking $w=\one_{\{S=N\}}$ shows
the minimum is at most $f'(N)=T-n\lambda$, so Conjecture~\ref{conj:floor} is
exactly the assertion
\[
 \sum_S w_Sf'(S)\ \ge\ f'(N)\qquad\text{for every fractional cover }w. \tag{2}
\]

\begin{proposition}[balanced covers]\label{prop:balanced}
Inequality (2) holds for every exact (balanced) cover, i.e.\ whenever
$\sum_{S\ni i}w_S=1$ for all $i$.
\end{proposition}

\begin{proof}
Balancedness gives $\sum_Sw_S|S|=\sum_i\sum_{S\ni i}w_S=n$, so, using
$f\ge g$ and the modularity of $g$,
\[
 \sum_Sw_Sf'(S)=\sum_Sw_Sf(S)-\lambda n
 \ \ge\ \sum_Sw_Sg(S)-\lambda n
 =\sum_i g_i\sum_{S\ni i}w_S-\lambda n
 =g(N)-\lambda n=f'(N),
\]
the last step because $g(N)=T$ and $f'(N)=T-\lambda n$.
The same computation works with $h$ in place of $g$.
\end{proof}

What is missing is the case of covers that strictly over-cover some players.
Writing $e_i:=\sum_{S\ni i}w_S-1\ge0$, the same computation gives
\[
 \sum_Sw_Sf'(S)\ \ge\ f'(N)
 +\max\Bigl(\sum_{i\in A}g_ie_i,\ \sum_{i\in B}h_ie_i\Bigr)
 -\lambda\sum_{i\in N}e_i,
\]
and since $\lambda\le f(\{i\})$, which equals $g_i$ for $i\in A$ and $h_i$
for $i\in B$, each inner sum dominates $\lambda$ times its \emph{own} block's
excess but not the other block's.  Closing that gap, or constructing the
attaining core allocation directly, is the remaining work.

\subsection{The leximax-min ceiling is not dual}

The mirror-image statement is an inequality but not an identity, and this is
the first place where the two selections part company.

\begin{observation}[$M$-side asymmetry]\label{obs:Mside}
For every core allocation, and in particular for $\Lmax(v)$,
\[
 M\ \ge\ \max_{\substack{P\subseteq A,\ Q\subseteq B\\ (P,Q)\ne(\emptyset,\emptyset)}}
 \frac{\min\bigl(g(P),h(Q)\bigr)}{|P|+|Q|}. \tag{3}
\]
Inequality (3) is not an identity: it fails strictly already at $n=6$.
\end{observation}

\begin{proof}[Proof of the inequality]
Let $x\in\core(v)$ and put $S=P\cup Q$.  Then $g(S)=g(P)$ and $h(S)=h(Q)$
because $g$ vanishes on $B$ and $h$ on $A$, so
\[
 \max_i x_i\ \ge\ \frac{x(S)}{|S|}\ \ge\ \frac{v(S)}{|S|}
 =\frac{\min\bigl(g(P),h(Q)\bigr)}{|P|+|Q|}.
\]
Apply this to $x=\Lmax(v)$.  Equivalently, (3) is Theorem~\ref{thm:duality}
read on complements.
\end{proof}

\paragraph{How to read (3).}
The maximum runs over pairs $(P,Q)$ other than $(\emptyset,\emptyset)$, with
the convention $g(\emptyset)=h(\emptyset)=0$.  There is no infinity
convention: a pair with an empty part contributes $\min(\cdot,0)=0$ and is
simply inactive, so only pairs with \emph{both} parts nonempty can bind.
This is visible in the proof, where an empty part makes $v(S)$ the minimum of
something and $0$.  Reading an empty part as $+\infty$ instead would reduce
(3) to a maximum over singletons and give a different --- and wrong --- value;
the misreading is easy enough that we spell the convention out.

\paragraph{The smallest counterexample.}
Scanning the complete two-generator enumeration in increasing $T$, no
counterexample to equality in (3) exists at $n\le5$ for any $T$, nor at $n=6$
for $T\le4$.  At $n=6$, $T=5$:
\[
 g=(3,2,0,0,0,0),\qquad h=(0,0,2,1,1,1),
\]
\[
 \Lmin(v)=\Bigl(1,\tfrac23,\tfrac43,\tfrac23,\tfrac23,\tfrac23\Bigr),
 \qquad
 \Lmax(v)=\Bigl(\tfrac65,\tfrac45,\tfrac65,\tfrac35,\tfrac35,\tfrac35\Bigr),
\]
so $M=6/5$, whereas the right-hand side of (3) equals $1$, attained for
instance at $P=\{1\},Q=\{3\}$ with $\min(3,2)/2=1$.  The gap is $1/5$.  Two
further counterexamples occur at $n=6$, $T=5$, with $g=(3,1,1)$, $h=(2,2,1)$
and with $g=(2,2,1)$, $h=(3,1,1)$ on the respective supports; all have
$M=6/5$ against a bound of $1$.  Over the same random samples used for
Conjecture~\ref{conj:floor}, (3) is tight in $120/120$ games at $n=4$,
$218/218$ at $n=5$, $160/165$ at $n=6$ and $74/83$ at $n=7$.

\begin{remark}[diagnosis]\label{rem:asymmetry}
Theorem~\ref{thm:floor} bounds the floor using only the upper caps
$x(S)\le f(S)$, and a single tightest cap can be saturated on its own, which
is why one coalition certifies the optimum.  Bounding the ceiling instead
requires the lower bounds $x(N\setminus S)\ge T-f(S)$, and \emph{minimizing}
a maximum must respect several of them simultaneously; no single one
certifies the optimum.  This is a genuine asymmetry between the two
selections, and it is the reason $\Lmax$ is the harder object throughout.
\end{remark}

\section{The master family}\label{sec:family}

\begin{definition}\label{def:family}
For integers $p\ge1$ and $k\ge2$ put $n=p+1+k$ and $T=pk$, and let $G(p,k)$
be the lower envelope of
\[
 g=\bigl(p(k-1),\ \underbrace{1,\dots,1}_{p},\ \underbrace{0,\dots,0}_{k}\bigr),
 \qquad
 h=\bigl(\underbrace{0,\dots,0}_{p+1},\ \underbrace{p,\dots,p}_{k}\bigr).
\]
We index the players $0,1,\dots,p+k$ and call player $0$ the \emph{head},
$O:=\{1,\dots,p\}$ the \emph{ones}, and $B:=\{p+1,\dots,p+k\}$; so
$A=\{0\}\cup O$ has $p+1$ elements and $|B|=k$.
\end{definition}

Both generators are nonnegative with total $p(k-1)+p=pk=T$, their supports
are disjoint, and every coordinate vanishes in one of them; by
Lemma~\ref{lem:envelope}, $G(p,k)$ is exact and singleton-normalized.

\begin{theorem}[closed forms]\label{thm:family}
For all $p\ge1$, $k\ge2$,
\begin{align*}
 \Lmin\bigl(G(p,k)\bigr)&=\Bigl(\tfrac{p^2(k-1)}{p+1},\ \bigl[\tfrac{p}{p+1}\bigr]^{p+k}\Bigr),\\
 \Lmax\bigl(G(p,k)\bigr)&=\Bigl(\tfrac{p(k-1)}{k},\ \bigl[\tfrac1k\bigr]^{p},\ \bigl[\tfrac{p(k-1)}{k}\bigr]^{k}\Bigr),\\
 \bigl\|\Lmin-\Lmax\bigr\|_1&=\frac{2p\bigl(p(k-1)-1\bigr)}{p+1},
 \qquad
 \rho\bigl(G(p,k)\bigr)=\frac{2\bigl(p(k-1)-1\bigr)}{k(p+1)},
\end{align*}
where the exponents indicate the number of coordinates, listed in the order
head, ones, $B$.
\end{theorem}

The proof occupies Sections~\ref{sec:keycap}--\ref{sec:gap}: a key cap, the
two optimizations, a membership check, and the arithmetic of the gap.

\subsection{The key cap}\label{sec:keycap}

For each $b\in B$,
\[
 f\bigl(O\cup\{b\}\bigr)=\max\bigl(g(O),h(\{b\})\bigr)=\max(p,p)=p,
\]
so by Theorem~\ref{thm:duality} every core allocation $x$ satisfies
\[
 x_1+\dots+x_p+x_b\ \le\ p\qquad\text{for every }b\in B. \tag{KC}
\]
That single family of $(p+1)$-player caps drives everything below.  Note that
the two generators \emph{tie} at these coalitions: the maximum defining $f$
is attained on both sides.

\subsection{Leximin}

Let $x\in\core(G(p,k))$ and $m:=\min_ix_i$.  Applying (KC) with any fixed $b$
gives $(p+1)m\le p$, so
\[
 m\le\frac{p}{p+1},\qquad\text{hence}\qquad \mu\le\frac{p}{p+1}.
\]
Suppose $m=p/(p+1)$.  Then in (KC) the $p+1$ numbers $x_1,\dots,x_p,x_b$ are
each at least $p/(p+1)$ and sum to at most $p=(p+1)\cdot p/(p+1)$, so all of
them equal $p/(p+1)$.  This holds for \emph{every} $b\in B$, hence
\[
 x_1=\dots=x_p=x_{p+1}=\dots=x_{p+k}=\frac{p}{p+1},
\]
and efficiency forces the head to
\[
 x_0=T-(p+k)\frac{p}{p+1}=pk-\frac{p(p+k)}{p+1}=\frac{p^2(k-1)}{p+1},
\]
using $k-(p+k)/(p+1)=(kp+k-p-k)/(p+1)=p(k-1)/(p+1)$.  Thus the max-min
allocation is \emph{unique}, and a unique max-min allocation is a fortiori
the leximin allocation.  Since $p(k-1)\ge1$, the head's coordinate is at
least $p/(p+1)$, so the floor of this allocation is indeed $p/(p+1)$.  It
remains to check that this point is in the core, which is done in
Section~\ref{sec:membership}.

\subsection{Leximax-min}

Let $x\in\core(G(p,k))$, $M:=\max_ix_i$ and $s:=x_1+\dots+x_p$.  By (KC),
$x_b\le p-s$ for every $b\in B$, and trivially $x_b\le M$ and $x_0\le M$.
Efficiency gives
\[
 pk=x_0+s+\sum_{b\in B}x_b\ \le\ M+s+k\min(M,\,p-s). \tag{4}
\]

\emph{Case A: $M\le p-s$.}  Then (4) reads $pk\le(k+1)M+s$.  The case
hypothesis is $s\le p-M$, so $pk\le(k+1)M+p-M=kM+p$, i.e.
$M\ge p(k-1)/k$.

\emph{Case B: $M>p-s$.}  Then (4) reads $pk\le M+s+k(p-s)$, i.e.
$(k-1)s\le M$, so $s\le M/(k-1)$.  The case hypothesis then gives
$M>p-s\ge p-M/(k-1)$, i.e.\ $Mk/(k-1)>p$, i.e.\ $M>p(k-1)/k$.

In both cases $M\ge p(k-1)/k$, so the leximax-min ceiling is at least
$p(k-1)/k$.  Equality can occur only in Case A with every inequality used
tight, that is with $x_0=M$, with $x_b=M$ for all $b\in B$, and with
\[
 s=pk-(k+1)M=pk-(k+1)\frac{p(k-1)}{k}=\frac{p\bigl(k^2-(k+1)(k-1)\bigr)}{k}=\frac pk .
\]
(Consistency: $p-s=p-p/k=p(k-1)/k=M$, so Case A holds with equality, as it
must.)  Thus at the min-max level the head and all of $B$ are pinned at
$p(k-1)/k$ and the ones share exactly $p/k$; since $p/k\le p(k-1)/k$ for
$k\ge2$, the ones occupy the last $p$ positions of the decreasing
rearrangement.  The next leximax stage minimizes their maximum subject to
their sum being $p/k$, whose unique optimum is $1/k$ each.  Hence $\Lmax$ is
as claimed, again modulo membership.

\subsection{Membership}\label{sec:membership}

This step is load-bearing: the two optimizations above determine the optima
only conditionally on the described allocations being feasible.

Index as in Definition~\ref{def:family} and, for $S\subseteq N$, put
\[
 \varepsilon:=[\,0\in S\,],\qquad i:=|S\cap O|,\qquad j:=|S\cap B|,
\]
so that $0\le i\le p$, $0\le j\le k$ and
\[
 g(S)=\varepsilon\,p(k-1)+i,\qquad h(S)=jp,\qquad
 f(S)=\max\bigl(\varepsilon\,p(k-1)+i,\ jp\bigr). \tag{5}
\]

\begin{lemma}\label{lem:membership}
Both allocations of Theorem~\ref{thm:family} lie in $\core(G(p,k))$ for all
$p\ge1$, $k\ge2$.  Moreover, for both of them the tight coalitions are
exactly
\[
\begin{array}{ll}
 \text{(T1)}\quad S=O\cup\{b\},\ b\in B, & f(S)=\max(p,p)=p,\\[2pt]
 \text{(T2)}\quad S=\{0\}\cup B',\ |B'|=k-1, & f(S)=\max\bigl(p(k-1),(k-1)p\bigr)=p(k-1),\\[2pt]
 \text{(T3)}\quad S=N, & f(N)=T.
\end{array}
\]
\end{lemma}

\begin{proof}
\emph{The leximin allocation.}  Write $q:=p/(p+1)$ and
$H:=p^2(k-1)/(p+1)=p(k-1)q$, so $x^*=(H,q,\dots,q)$ with $p+k$ copies of $q$.
Efficiency:
\[
 H+(p+k)q=q\bigl(p(k-1)+p+k\bigr)=q(pk+k)=k(p+1)\frac{p}{p+1}=pk=T,
\]
and $x^*>0$.  It remains to check $x^*(S)=\varepsilon H+(i+j)q\le f(S)$.

$\varepsilon=0$, $j=0$: $S$ is a subset of $O$, and $x^*(S)=iq<i=f(S)$ for
$i\ge1$ because $q<1$; for $i=0$, $S=\emptyset$.  Strict.

$\varepsilon=0$, $j\ge1$: bound $f(S)\ge jp$ and use $i\le p$:
\[
 x^*(S)=(i+j)\frac{p}{p+1}\le(p+j)\frac{p}{p+1}\le jp
 \iff p+j\le j(p+1)\iff p\le jp,
\]
which holds since $j\ge1$.  Equality throughout requires $i=p$ and $j=1$,
i.e.\ $S=O\cup\{b\}$, which is (T1).

$\varepsilon=1$, $j\le k-1$: put $G_0:=p(k-1)+i=g(S)$, so $x^*(S)=q(G_0+j)$
and
\[
 x^*(S)\le G_0\iff p(G_0+j)\le(p+1)G_0\iff jp\le G_0=p(k-1)+i,
\]
which holds because $jp\le(k-1)p$.  Hence $x^*(S)\le g(S)\le f(S)$, with
equality exactly when $j=k-1$ and $i=0$, i.e.\ $S=\{0\}\cup B'$ with
$|B'|=k-1$, which is (T2).

$\varepsilon=1$, $j=k$: then $h(S)=kp=T$, so $f(S)\ge T\ge x^*(S)$, with
equality only for $S=N$ since $x^*>0$; this is (T3).

\emph{The leximax-min allocation.}  Write $L:=p(k-1)/k$ and $r:=1/k$, so
$y^*=(L,r,\dots,r,L,\dots,L)$ with $p$ copies of $r$ and $k$ copies of $L$.
Efficiency:
\[
 L+pr+kL=(k+1)L+\frac pk=\frac{p\bigl[(k+1)(k-1)+1\bigr]}{k}=\frac{pk^2}{k}=pk=T,
\]
and $y^*>0$.  Now $y^*(S)=\varepsilon L+ir+jL$.

$\varepsilon=0$, $j=0$: $y^*(S)=i/k\le i=f(S)$ since $k\ge2$, strictly for
$i\ge1$.

$\varepsilon=0$, $j\ge1$: comparing with the $h$-side,
\[
 \frac{i+jp(k-1)}{k}\le jp\iff i+jp(k-1)\le jpk\iff i\le jp,
\]
which holds since $i\le p\le jp$, with equality exactly when $i=p$ and $j=1$,
i.e.\ (T1).

$\varepsilon=1$, $j\le k-1$: $y^*(S)=(j+1)L+i/k$, and multiplying by $k>0$,
\[
 y^*(S)\le p(k-1)+i
 \iff (j+1)p(k-1)+i\le kp(k-1)+ki
 \iff p(j+1-k)\le i,
\]
after dividing by $k-1\ge1$; the left-hand side is at most $0\le i$, with
equality exactly when $j=k-1$ and $i=0$, i.e.\ (T2).

$\varepsilon=1$, $j=k$: as before $f(S)\ge h(S)=kp=T\ge y^*(S)$, with
equality only at $S=N$.
\end{proof}

\subsection{The hinge: both selections saturate the same caps}

Lemma~\ref{lem:membership} contains the expository point of the paper.  The
active constraints are the same for $\Lmin$ and $\Lmax$: the two families
(T1) and (T2), plus efficiency.  Both are coalitions at which the two
generators \emph{tie}, which is the structural reason they are the active
ones --- there is no slack on either side to give.  What differs is only how
the two selections distribute mass \emph{inside} a tight coalition.  On
$S=O\cup\{b\}$, which has total exactly $p$:
\[
\begin{array}{lcc}
 & \text{each of the }p\text{ ones} & \text{the single }b\in B\\[2pt]
 \Lmin: & \dfrac{p}{p+1} & \dfrac{p}{p+1}\\[10pt]
 \Lmax: & \dfrac1k & \dfrac{p(k-1)}{k}
\end{array}
\]
Leximin splits the budget $p$ evenly over all $p+1$ members; leximax-min
gives the ones the least it can and puts the whole remainder on the one
member of $B$, because loading $B$ is how it lowers the global maximum.  The
divergence is visible in this one coalition.  Everything else is bookkeeping:
leximin's even split leaves the residual $T-(p+k)p/(p+1)$ with nowhere to go
but the head, while leximax-min's loading of every $b\in B$ up to
$p(k-1)/k$ consumes almost all of $T$ across $B$.

\subsection{The gap}\label{sec:gap}

\begin{proof}[Proof of Theorem~\ref{thm:family}, conclusion]
Put $D:=p(k-1)-1\ge0$.  Coordinatewise,
\begin{align*}
 \text{head:}&& \frac{p^2(k-1)}{p+1}-\frac{p(k-1)}{k}&=\frac{p(k-1)D}{k(p+1)},\\
 \text{each of the $p$ ones:}&& \frac{p}{p+1}-\frac1k&=\frac{D}{k(p+1)},\\
 \text{each of the $k$ members of $B$:}&& \frac{p(k-1)}{k}-\frac{p}{p+1}&=\frac{pD}{k(p+1)},
\end{align*}
all three nonnegative.  Hence
\[
 \|\Lmin-\Lmax\|_1=\frac{p(k-1)D+pD+kpD}{k(p+1)}
 =\frac{D\bigl(p(k-1)+p+kp\bigr)}{k(p+1)}
 =\frac{2pkD}{k(p+1)}=\frac{2pD}{p+1},
\]
and $\rho=2pD/\bigl((p+1)pk\bigr)=2D/\bigl(k(p+1)\bigr)$.
\end{proof}

\begin{corollary}[deficit identity]\label{cor:deficit}
$\displaystyle 2-\rho\bigl(G(p,k)\bigr)=\frac{2n}{k(p+1)}$, where $n=p+1+k$.
\end{corollary}

\begin{proof}
$2-\dfrac{2\bigl(p(k-1)-1\bigr)}{k(p+1)}
 =\dfrac{2\bigl[k(p+1)-p(k-1)+1\bigr]}{k(p+1)}
 =\dfrac{2(kp+k-pk+p+1)}{k(p+1)}=\dfrac{2(k+p+1)}{k(p+1)}$.
\end{proof}

\begin{remark}[the deficit is twice the overlap]\label{rem:overlap}
Corollary~\ref{cor:deficit} also follows from Lemma~\ref{lem:ceiling}.  In
$G(p,k)$ the coordinatewise minima of the two allocations are $p(k-1)/k$ at
the head, $1/k$ at each one, and $p/(p+1)$ at each member of $B$, so
\[
 \sum_i\min\bigl(\Lmin_i,\Lmax_i\bigr)
 =\frac{p(k-1)}{k}+\frac pk+\frac{kp}{p+1}=p+\frac{kp}{p+1}<p+k=n-1,
\]
and $2\cdot\bigl(p+kp/(p+1)\bigr)/(pk)=2(p+1+k)/\bigl(k(p+1)\bigr)$.  The
overlap between the two selections is thus bounded by $n$ while $T=pk$ grows
quadratically: this is the precise sense in which their supports become
asymptotically disjoint.
\end{remark}

\subsection{The supremum}

\begin{theorem}[value of the supremum]\label{thm:sup}
For every $n\ge4$,
\[
 s(n)\ \ge\ 2-\frac{8n}{n^2-(n\bmod2)}
 =\begin{cases}
   2-\dfrac8n,& n\text{ even},\\[6pt]
   2-\dfrac{8n}{n^2-1},& n\text{ odd},
  \end{cases}
\]
attained on the two-generator class by $G(p,k)$ with the balanced split
$k=\lceil n/2\rceil$, $p+1=\lfloor n/2\rfloor$.  Consequently
$\sup_n s(n)=2$.
\end{theorem}

\begin{proof}
By Corollary~\ref{cor:deficit} the deficit of $G(p,k)$ is $2n/(k(p+1))$ with
$k+(p+1)=n$, $k\ge2$ and $p+1\ge2$.  Minimizing the deficit means maximizing
the product $k(p+1)$ of two positive integers with fixed sum $n$, whose
maximum is $\lfloor n/2\rfloor\lceil n/2\rceil$, equal to $n^2/4$ for even
$n$ and $(n^2-1)/4$ for odd $n$, attained at the balanced split (for $n\ge4$
this split satisfies the constraints $k\ge2$, $p\ge1$).  Substituting gives
the displayed bound.  For the last claim, $\rho\le2$ always by
Lemma~\ref{lem:ceiling}, and $2-8/n\to2$.
\end{proof}

Table~\ref{tab:values} lists the resulting values.  The ratio crosses $1$ at
$n=8$: from eight players on, the two egalitarian core selections can differ
by more than the entire worth of the grand coalition.

\begin{table}[t]
\centering
\begin{tabular}{@{}rlll@{}}
\toprule
$n$ & $(p,k)$ & lower bound on $s(n)$ & decimal\\
\midrule
 6 & $(2,3)$ & $2/3$ & $0.6667$\\
 7 & $(2,4)$ & $5/6$ & $0.8333$\\
 8 & $(3,4)$ & $1$ & $1.0000$\\
 9 & $(3,5)$ & $11/10$ & $1.1000$\\
10 & $(4,5)$ & $6/5$ & $1.2000$\\
11 & $(4,6)$ & $19/15$ & $1.2667$\\
12 & $(5,6)$ & $4/3$ & $1.3333$\\
16 & $(7,8)$ & $3/2$ & $1.5000$\\
20 & $(9,10)$ & $8/5$ & $1.6000$\\
100 & $(49,50)$ & $48/25$ & $1.9200$\\
\bottomrule
\end{tabular}
\caption{The bound of Theorem~\ref{thm:sup} at the balanced split.  All
entries are exact; decimals are rounded.}
\label{tab:values}
\end{table}

\begin{remark}[the degenerate corner]\label{rem:degenerate}
$G(1,2)$ has $n=4$, $T=2$, $g=(1,1,0,0)$, $h=(0,0,1,1)$ and
$\rho=2(1\cdot1-1)/(2\cdot2)=0$, consistent with the four-player rigidity;
the deficit identity reads $2-0=2n/(k(p+1))=8/4=2$.
\end{remark}

\begin{remark}[there is nothing special about $|B|=4$]\label{rem:notfour}
An earlier version of this family fixed $k=|B|=4$, giving
$\rho=(3p-1)/\bigl(2(p+1)\bigr)\to3/2$.  Corollary~\ref{cor:deficit} shows
that $k$ is a free parameter and that the optimum is $k\approx n/2$, so the
limit $3/2$ was an artifact of the fixed block size; the $k=4$ slice is
already beaten from $n=10$ on ($7/6$ against $6/5$).  The number $4$ that
appears as $|\supp(\Lmax)|$ in the $n=5$ extremal class
(Section~\ref{sec:fixedn}) is a small-$n$ phenomenon and not a feature of the
asymptotics: in $G(p,k)$ every coordinate of $\Lmax$ is strictly positive, so
$|\supp(\Lmax)|=n$ and no player is starved at all.
\end{remark}

\section{Fixed $n$: enumeration, certificates, and the $n=5$ anomaly}
\label{sec:fixedn}

\subsection{Complete two-generator enumeration}

Write $s_2(n)$ for the supremum of $\rho$ over singleton-normalized
two-generator exact games on $n$ players.  Table~\ref{tab:enum} reports
complete enumerations, up to player permutation, generator swap and scaling,
of all integer generator pairs with common total in the stated range.  All
arithmetic is exact rational.

\begin{table}[t]
\centering
\begin{tabular}{@{}rrrll@{}}
\toprule
$n$ & totals & games & enumerated max $\rho$ & Theorem~\ref{thm:sup}\\
\midrule
3 & $2$--$30$ & $225$ & $0$ & ---\\
4 & $2$--$30$ & $3{,}040$ & $0$ & $0$\\
5 & $2$--$24$ & $4{,}357$ & $1/2$ & $1/3$\\
6 & $2$--$20$ & $8{,}213$ & $2/3$ & $2/3$\\
7 & $2$--$14$ & $1{,}609$ & $5/6$ & $5/6$\\
8 & $2$--$12$ & $1{,}049$ & $1$ & $1$\\
\bottomrule
\end{tabular}
\caption{Complete two-generator enumerations.  The rows $n=3,4$ are controls:
they reproduce the four-player rigidity.  The last column is the value of the
balanced-split family bound.}
\label{tab:enum}
\end{table}

Two readings must be kept apart.

\paragraph{What is complete.}
Because $\rho$ is scale-invariant, every two-generator game with rational
generators has a unique primitive integer representative, and enumeration to
total $T_{\max}$ is exactly complete for the games whose primitive
representative has total at most $T_{\max}$ --- equivalently, for rational
shapes of denominator at most $T_{\max}$.  It says nothing about finer
denominators or about irrational generators, so enumeration-to-$T$ is a
lower-bound method for $s_2(n)$: the entries of Table~\ref{tab:enum} are
maxima over the enumerated range, not values of $s_2$.  As a partial probe of
the excluded region, shape hill-climbing at fine resolution was run from the
known witnesses (resolution $1/60$ at $n=5$; $1/56$ and $1/60$ at $n=6$;
$1/64$ at $n=7$).  None improved on the tabulated maximum; at $n=7$ the climb
from the $5/6$ witness stalled after $42$ steps, and a second start from a
base of $2/3$ climbed only to $3/4$.

\paragraph{What agrees.}
The enumerated maxima agree exactly with Theorem~\ref{thm:sup} at $n=6,7,8$,
so on the two-generator class that bound is attained there, and
\[
 s_2(n)\ \stackrel{?}{=}\ 2-\frac{8n}{n^2-(n\bmod2)}\qquad(n\ge6)
\]
is a live candidate, currently confirmed at three values of $n$.  At $n=5$
the two disagree, and the enumeration is the correct one; see
Section~\ref{sec:n5}.

\subsection{Searching the full exact class}

The two-generator restriction is not innocuous.  Adding a generator
\emph{lowers} the envelope and hence \emph{grows} the core, so it is not
automatic that two generators are extremal.  The following observation lets
one search the entire exact class rather than a bounded-generator slice.

\begin{lemma}[core closing]\label{lem:closing}
Fix $T>0$ and a vector $c=(c_S)$ indexed by the proper coalitions, and put
$P(c)=\{x: x(N)=T,\ x(S)\ge c_S\ \text{for all proper }S\}$.  If
$P(c)\ne\emptyset$, then the game $v$ defined by $v(S)=\min_{x\in P(c)}x(S)$
for proper $S$ and $v(N)=T$ is exact, and $\core(v)=P(c)$.
\end{lemma}

\begin{proof}
Every $x\in P(c)$ satisfies $x(S)\ge v(S)$ by definition of $v$ and
$x(N)=T$, so $P(c)\subseteq\core(v)$; and each minimum is attained, so $v$ is
exact.  Conversely, $v(S)\ge c_S$ for every proper $S$ because every point of
$P(c)$ satisfies $x(S)\ge c_S$; hence $x\in\core(v)$ implies $x(S)\ge
v(S)\ge c_S$ and $x(N)=T$, i.e.\ $x\in P(c)$.
\end{proof}

Perturbing $c$ and re-closing therefore moves within the exact class and
reaches cores with arbitrarily many vertices.  Hill climbing in this
``$c$-space'', \emph{seeded at the extremal two-generator game} of
Table~\ref{tab:enum}, produced no improvement at $n=5$ ($240$ steps, stays at
$1/2$), $n=6$ ($496$ steps, stays at $2/3$) or $n=7$ ($301$ steps, stays at
$5/6$).  Since the search is over all exact games, this certifies the
two-generator optimum as a \emph{local} optimum in the full exact class at
those three values of $n$.  Independently, from the known extremal at each of
$n=5,6,7,8$, about $200$ structured and random third generators at fixed
common total were added: zero improvements.

One limitation must be stated.  Unseeded $c$-space restarts reach only $1/3$
at $n=6$, far below the known $2/3$; that variant of the search is
under-powered for $n\ge6$ and certifies nothing.  Only the seeded runs
support the local-optimality claim, and none of these runs bears on global
optimality.

\subsection{The $n=5$ anomaly}\label{sec:n5}

At $n=5$ the balanced-split family gives only $1/3$, but the truth is $1/2$,
and the witness has a non-uniform $h$-block.  The proof is by hand from the
four core vertices; no linear program is used.

\begin{theorem}[$n=5$]\label{thm:n5}
Let $w$ be the lower envelope of $g^1=(1,0,0,1,2)$ and $g^2=(0,1,3,0,0)$.
Then $w$ is exact and singleton-normalized with $w(N)=4$, and
\[
 \Lmin(w)=\Bigl(\tfrac12,\tfrac12,\tfrac54,\tfrac12,\tfrac54\Bigr),\qquad
 \Lmax(w)=(1,0,1,1,1),\qquad
 \|\Lmin-\Lmax\|_1=2,\qquad \rho(w)=\tfrac12 .
\]
In particular $s(5)\ge1/2>1/3$.
\end{theorem}

\begin{proof}
Exactness and normalization are immediate from Lemma~\ref{lem:envelope}, each
coordinate vanishing in one generator.  Exact rational vertex enumeration
gives
\[
 \core(w)=\operatorname{conv}\{V_1,V_2,V_3,V_4\},
\]
\[
 V_1=(0,1,2,0,1),\quad V_2=(0,1,3,0,0),\quad V_3=(1,0,0,1,2),\quad V_4=(1,0,1,1,1).
\]
Write $x=aV_1+bV_2+cV_3+dV_4$ with $a+b+c+d=1$, $a,b,c,d\ge0$, and put
$s:=c+d\in[0,1]$.  Then
\[
 x_1=x_4=s,\qquad x_2=1-s,\qquad x_3+x_5=3-s. \tag{6}
\]

\emph{Leximin.}  By (6), $\min_ix_i\le\min(s,1-s)\le1/2$, with equality
forcing $s=1/2$ and hence $x_1=x_2=x_4=1/2$.  Then $x_3+x_5=5/2$, and the
leximin continuation splits this evenly, giving $5/4$ each.  The resulting
allocation is attained: $(a,b,c,d)=(1/4,1/4,1/2,0)$ gives $s=1/2$ and
$x=(1/2,1/2,5/4,1/2,5/4)$ exactly.  Hence $\Lmin(w)$ is as claimed.

\emph{Leximax-min.}  By (6), $\max_ix_i\ge(x_3+x_5)/2=(3-s)/2\ge1$ since
$s\le1$.  Equality forces $s=1$, hence $a=b=0$, and $x_3=x_5=1$.  Among
$c+d=1$ the condition $x_3=d=1$ forces $d=1$, $c=0$, i.e.\ $x=V_4$, uniquely.
Hence $\Lmax(w)=(1,0,1,1,1)$.

Finally $\|\Lmin-\Lmax\|_1=\tfrac12+\tfrac12+\tfrac14+\tfrac12+\tfrac14=2$
and $\rho=2/4=1/2$.
\end{proof}

Four independent methods agree on this game: the hand argument above, exact
rational core-vertex enumeration, a staged exact solver over linear real
arithmetic, and a from-scratch exact-rational two-phase simplex with Bland's
rule.

\begin{conjecture}[$n=5$ supremum]\label{conj:n5}
$s(5)=1/2$.
\end{conjecture}

\paragraph{Evidence.}
Roughly $4\times10^4$ distinct exact five-player games were evaluated in
exact rational arithmetic with nothing above $1/2$: the complete
two-generator enumeration of Table~\ref{tab:enum} ($4{,}357$ games, maximum
$1/2$); complete three-generator enumerations at totals $4,5,6,7$ ($335$,
$1{,}175$, $3{,}747$ and $9{,}923$ games, maxima $1/2$, $2/5$, $1/3$,
$8/21$); $c$-space hill climbing in the sense of Lemma~\ref{lem:closing}
($68$ restarts, about $6{,}600$ evaluations), which plateaus at $1/2$; and
fine-denominator shape climbing from the extremal witness scaled to total
$60$, which does not improve.  The controls at $n=3,4$ in
Table~\ref{tab:enum} return maximum $0$ from the same pipeline.

\paragraph{Structure of the extremal class.}
The following are computational findings over $360$ distinct profiled games
at $n=5$, of which $199$ have $\rho>1/3$ and $107$ attain $\rho=1/2$.  Among
the $107$ extremal games, without exception:
\begin{itemize}[nosep]
\item $\Lmin$ has stratum composition $(3,2)$, and $\Lmax$ has stratum
      composition $(1,4)$;
\item $\Lmax$ has exactly one zero coordinate, so $|\supp(\Lmax)|=4$;
\item normalized to $T=4$ the pair is unique up to permutation:
      $\Lmin\sim(\tfrac12,\tfrac12,\tfrac12,\tfrac54,\tfrac54)$ and
      $\Lmax\sim(0,1,1,1,1)$ --- that is, exactly the pair of
      Theorem~\ref{thm:n5}.
\end{itemize}
Among the $199$ games above $1/3$, the $\Lmin$ strata are $(3,2)$ in $196$
cases and $(2,1,2)$ in $3$; the pattern $(4,1)$ never occurs above $1/3$, and
$\Lmax$ is never the uniform allocation.

One tempting invariant is refuted.  Call $S$ a \emph{flat split} if
$v(S)+v(N\setminus S)=v(N)$; then $x(S)$ is constant on the core, which is
the affine seesaw visible in (6) --- in the game of Theorem~\ref{thm:n5},
$w(\{1,2\})+w(\{3,4,5\})=1+3=4$.  Only $76$ of the $107$ extremal games have
any flat split, and $31$ attain exactly $1/2$ with none, so the seesaw is a
feature of the particular presentation in Theorem~\ref{thm:n5} and not an
invariant of extremality.  When a flat split does exist its sizes are always
$\{2,3\}$, and the $\Lmax$-starved player lies on the size-$2$ side in all
$76$ cases.

\paragraph{Two regimes.}
Theorem~\ref{thm:sup} gives only $1/3$ at $n=5$ because the family $G(p,k)$
has a uniform $h$-block and the available splits $(k,p+1)\in\{(2,3),(3,2)\}$
are too coarse; the extremal five-player game has a non-uniform $h$-block.
The crossover to the $G$-family shape happens at $n=6$.  The two regimes
differ qualitatively and not just numerically: at $n=5$ the extremal $\Lmax$
starves one player to $0$ and has support of size $4$, whereas in $G(p,k)$
the extremal $\Lmax$ is strictly positive everywhere.  Explaining the
crossover is open (Section~\ref{sec:scope}).


\subsection{One-coalition certificates and the regime change}
\label{sec:n5closed}

Write, for any singleton-normalized exact game with the general dual bound
$f(S)=T-v(N\setminus S)$,
\[
 \lambda:=\min_{\emptyset\ne S}\frac{f(S)}{|S|},\qquad
 \nu:=\max_{\emptyset\ne S}\frac{v(S)}{|S|},\qquad
 \mu:=\min_i\Lmin(v)_i,\qquad M:=\max_i\Lmax(v)_i .
\]
The general duality gives the bracketing $\mu\le\lambda$ and $M\ge\nu$
at once.  Observation~\ref{obs:Mside} says the second bracket can be
strict from $n=6$ on.  Two facts sharpen this into a statement about the
regime change at five players.

\begin{proposition}[covering obstruction]\label{prop:cover}
Let $\tau_\nu:=\{S:v(S)=\nu|S|\}$ be the $\nu$-tight family.  If $\tau_\nu$
covers $N$ and $\nu>T/n$, then no core point satisfies $x\le\nu\one$, and
hence $M>\nu$.
\end{proposition}

\begin{proof}
Sum the core inequalities $x(S)\ge v(S)=\nu|S|$ over a cover with weights
$w\ge0$: with $c_i:=\sum_{S\ni i}w_S\ge1$ this gives
$\sum_ic_i(x_i-\nu)\ge0$.  If $x\le\nu\one$, every term is nonpositive, so
$x_i=\nu$ for every $i$; then $T=\nu n$, contradicting $\nu>T/n$.
\end{proof}

Moreover no \emph{balanced} collection can witness $M>\nu$: for balanced $w$,
$\sum_Sw_Sv(S)\le T$ because the core is nonempty, and the same computation
then yields $\sum_Sw_S\bigl(\nu|S|-v(S)\bigr)\ge\nu n-T$, which is what
equality requires.  So a failure of $M=\nu$ needs a family of $\nu$-tight
coalitions that covers $N$ but contains no balanced subcollection.  At the
smallest counterexample of Observation~\ref{obs:Mside} exactly such a
family occurs: the four $\nu$-tight sets $\{2,3\},\{1,3,4\},\{1,3,5\},
\{1,3,6\}$ cover $N$ with per-player excess $(2,0,3,0,0,0)$, and player $2$,
lying only in $\{2,3\}$, blocks every balanced subfamily.  At the
five-player extremal game of Theorem~\ref{thm:n5}, by contrast, the optimal
cover of the dual program is balanced and attains equality --- the
obstruction is absent exactly where the conjecture asserts it must be.  The
covering test is sufficient for $M>\nu$ but not necessary: at six players,
games with $M>\nu$ outnumber those with a covering $\nu$-tight family.

\begin{conjecture}[one-coalition certificates at five players]
\label{conj:n5closed}
For every singleton-normalized exact game on five players, $\mu=\lambda$ and
$M=\nu$: both egalitarian levels are certified by a single coalition.
\end{conjecture}

The conjecture is exact in $10{,}813$ five-player games, of which $9{,}614$
belong to complete enumerations, with zero violations; and the $M$-side is
provably false from six players on (Observation~\ref{obs:Mside}).  If true,
it is the formal content of the regime change of
Section~\ref{sec:n5}: at five players both selections still admit
one-coalition certificates, and $s(5)=\tfrac12$ is governed by a single
bottleneck coalition, while from six players the ceiling requires genuinely
fractional covers --- which is precisely where $s(n)$ begins its climb
toward~$2$.

\section{Obstructions to an upper bound}\label{sec:obstruction}

\subsection{No constant below $2$}

Theorem~\ref{thm:sup} already rules out one shape of answer.  Any bound of
the form $\|\Lmin(v)-\Lmax(v)\|_1\le c\cdot v(N)$ with a constant $c<2$,
valid for all $n$, is false: choose $n$ even with $2-8/n>c$.  So the only
possible clean upper bound is dimension-dependent, and the right target is
the finite-$n$ statement
\[
 s_2(n)=2-\frac{8n}{n^2-(n\bmod2)}\qquad(n\ge6),
\]
currently confirmed by complete enumeration only at $n=6,7,8$, together with
the question of whether $s(n)=s_2(n)$.

\subsection{No bound linear in the leximin floor}

A second shape is also excluded, and this one was a live candidate.  The
inequality
\[
 \|\Lmin(v)-\Lmax(v)\|_1\ \le\ 4(n-4)\,\mu(v),\qquad \mu(v)=\min_i\Lmin(v)_i,
\]
holds with no violations in $900$ fresh random five-player games and is tight
both on the $n=5$ extremal class and on $G(3,4)$ at $n=8$, where both sides
equal $12$.  It fails from $n=9$ on:
\begin{center}
\begin{tabular}{@{}llll@{}}
\toprule
game & $n$ & $\|\Lmin-\Lmax\|_1$ & $4(n-4)\mu$\\
\midrule
$G(4,4)$ & 9 & $88/5$ & $16$\\
$G(3,5)$ & 9 & $33/2$ & $15$\\
$G(4,5)$ & 10 & $24$ & $96/5$\\
$G(5,5)$ & 11 & $95/3$ & $70/3$\\
\bottomrule
\end{tabular}
\end{center}
The reason is structural rather than accidental: in $G(p,k)$ one has
$\mu=p/(p+1)<1$, bounded, while the gap grows like $2T$.  No bound of the
form $C(n)\mu(v)$ can track the divergence, whatever $C$ is.

\subsection{Non-monotonicity under majorization}

The most natural route to an upper bound on the two-generator class is to
argue that a uniform $h$-block is without loss of generality --- the extremal
family $G(p,k)$ has one --- and then to optimize over $g$ alone using the
floor formula of Conjecture~\ref{conj:floor}.  That route is closed.

Define a \emph{smoothing step} as a Robin Hood transfer of one unit from a
larger to a smaller $h$-coordinate, so that $h$ moves down in the
majorization order, toward uniform.  The premise that a smoothing step never
decreases $\rho$ is false, and not only for small blocks:
\begin{center}
\begin{tabular}{@{}lrrrr@{}}
\toprule
$|B|$ & $2$ & $3$ & $4$ & $5$\\
\midrule
smoothing steps tested & $40$ & $76$ & $84$ & $84$\\
violations & $18$ & $27$ & $11$ & $21$\\
rate & $45\%$ & $36\%$ & $13\%$ & $25\%$\\
\bottomrule
\end{tabular}
\end{center}
In particular it fails for $|B|\ge4$, so the block-size threshold one might
guess from the $n=5$ anomaly does not exist.  An explicit verified
counterexample at $|B|=4$, with $n=7$ and $T=10$, both allocations confirmed
in the core:
\[
 g=(0,5,0,3,2,0,0),\qquad h=(1,0,2,0,0,5,2),
\]
\[
 \Lmin=(1,2,1,1,1,3,1),\qquad
 \Lmax=\Bigl(\tfrac56,\tfrac52,\tfrac56,\tfrac43,\tfrac76,\tfrac52,\tfrac56\Bigr),
 \qquad \rho=\tfrac15 .
\]
Smoothing the $h$-block $(5,2)\mapsto(4,3)$, i.e.\ replacing $h$ by
$(1,0,2,0,0,4,3)$, gives
\[
 \Lmin=\Lmax=(1,2,1,1,1,2,2),\qquad \rho=0 .
\]
So a single Robin Hood step inside one generator collapses the ratio from
$1/5$ to $0$, and $\rho$ is not monotone under majorization of $h$.

The two facts sit together uncomfortably and both are true: smoothing makes a
game more symmetric and tends to \emph{collapse} the two selections onto each
other, yet the extremal family has a perfectly uniform $h$.  The dependence
is therefore genuinely non-monotone, and any upper-bound proof must handle
general $h$ directly rather than reducing to the uniform case.

\section{What the results do and do not establish}\label{sec:scope}

\paragraph{Proved.}
Lemma~\ref{lem:envelope} (lower-envelope generation and its converse);
Lemma~\ref{lem:ceiling} (ceiling and overlap identity);
Lemma~\ref{lem:disjoint} (support disjointness); Theorem~\ref{thm:duality}
(two-generator core duality, with identity (1) valid without disjointness);
Theorem~\ref{thm:floor} ($\mu\le\lambda$); Proposition~\ref{prop:balanced}
(the cover inequality (2) for balanced covers); the validity half of
Observation~\ref{obs:Mside}; Theorem~\ref{thm:family} with
Lemma~\ref{lem:membership} (closed forms for $G(p,k)$, including membership
for all $p\ge1$, $k\ge2$ and the exact list of tight coalitions);
Corollary~\ref{cor:deficit} and Remark~\ref{rem:overlap} (the deficit
identity, twice);  Theorem~\ref{thm:sup} ($s(n)\ge2-8n/(n^2-(n\bmod2))$ and
$\sup_ns(n)=2$); Lemma~\ref{lem:closing} (core closing);
Theorem~\ref{thm:n5} ($\rho=1/2$ at $n=5$, by hand).  The failure of the
$4(n-4)\mu$ bound from $n=9$, and the non-monotonicity of $\rho$ under
majorization of $h$, are proved by the explicit games displayed in
Section~\ref{sec:obstruction}: each is a finite exact-rational computation on
a named game, and each is a refutation, so a single instance suffices.

\paragraph{Enumerated, not proved.}
The entries of Table~\ref{tab:enum} are maxima over the stated integer
ranges, hence complete only for rational shapes whose primitive integer
representative has total in that range; they are lower bounds for $s_2(n)$,
not its value.  The strict failure of equality in (3) is established by the
displayed six-player game, but the claim that this is the \emph{smallest}
such failure rests on the enumeration.  The tightness rates for (3) and for
Conjecture~\ref{conj:floor} are sample statistics over $586$ random games.
The structural findings about the $n=5$ extremal class in
Section~\ref{sec:n5} --- unique solution pair up to permutation, $(3,2)$ and
$(1,4)$ stratum compositions, $|\supp(\Lmax)|=4$, the refutation of the
seesaw hypothesis --- are computational, over $360$ profiled games, and none
of them is proved.  The local-optimality certificates from seeded $c$-space
climbing at $n=5,6,7$ are numerical searches over the full exact class; the
unseeded variant is under-powered at $n\ge6$ and certifies nothing.

\paragraph{Conjectured.}
Five problems remain open.
\begin{enumerate}[label=(\arabic*),nosep]
\item \textbf{Conjecture~\ref{conj:floor}}, equality in the leximin floor
      formula.  Reduced by linear-programming duality to the fractional-cover
      inequality (2); proved for balanced covers
      (Proposition~\ref{prop:balanced}); the over-covering case is open.
      Verified in $586/586$ random two-generator games at $n=4,\dots,7$.
\item \textbf{An upper bound.}  Theorem~\ref{thm:sup} forces any true bound
      to be asymptotically $2$, so every bound $c\cdot v(N)$ with $c<2$ is
      false for large $n$.  The right target is
      $s_2(n)=2-8n/(n^2-(n\bmod2))$ for $n\ge6$, confirmed by complete
      enumeration at $n=6,7,8$ only.  Section~\ref{sec:obstruction} closes
      the majorization route and the $C(n)\mu$ route.
\item \textbf{Is the two-generator class extremal?}  Adding a generator
      lowers the envelope and grows the core, so this is not automatic.  All
      evidence is negative for a counterexample: zero improvements from about
      $200$ third generators at each of $n=5,6,7,8$, and no improvement from
      seeded $c$-space climbing over the full exact class at $n=5,6,7$.  Open
      globally.
\item \textbf{Conjecture~\ref{conj:n5}}, $s(5)=1/2$, together with the
      two-regime picture suggested by the extremal-class structure: $(3,2)$
      $\Lmin$-strata above $1/3$, and $(4,1)$ never exceeding $1/3$.
\item \textbf{A structural explanation of the $n=5$ versus $n\ge6$ regime
      change}: a starved $\Lmax$ with $|\supp(\Lmax)|=4$ at $n=5$ against a
      strictly positive $\Lmax$ in $G(p,k)$.
\end{enumerate}

\paragraph{Not claimed.}
Exactness theory, lower-envelope exactification, and the characterizations of
the egalitarian selections are prior art
\citep{ArinInarra2001,ArinKuipersVermeulen2003,Biswas1999,Dietzenbacher2021,
LlerenaMauri2017,CsokaHeringsKoczy2011,StudenyKratochvil2022}; the
four-player rigidity is the companion paper's classification
\citep{Flanagan2026} combined with the Lorenz implication of
\citet{LlerenaMauri2017}.  The candidate
new content is the two-generator core duality and its one-sided
floor certificate, the family $G(p,k)$ with its closed forms and deficit
identity, the resulting evaluation $\sup_n s(n)=2$, the identification of the
$n=5$ anomaly with a proof by hand, and the two closed routes of
Section~\ref{sec:obstruction}.  Novelty remains conditional on specialist
review.

\section*{Reproducibility note}

Every theorem above has a hand proof, and no statement labelled ``proved''
depends on a program.  Exact rational programs in the companion repository
independently reconstruct all displayed games and coalition values, the
closed forms of Theorem~\ref{thm:family} at twelve parameter pairs, the full
coalition-inequality check and the tight-set prediction of
Lemma~\ref{lem:membership} for every $(p,k)$ with $1\le p\le7$, $2\le k\le7$
and $n\le13$, and the four-way agreement on the game of
Theorem~\ref{thm:n5}.  They are bug audits, not proof dependencies.  All
arithmetic is exact rational; no floating-point value appears in any stated
result.  Reported optima come from a from-scratch exact-rational two-phase
simplex with Bland's rule, from a staged exact solver over linear real
arithmetic, or from exact vertex enumeration, with at least two methods
agreeing on every displayed allocation; one library simplex implementation
cycles on these degenerate programs and was not used for any reported number.

\begin{thebibliography}{99}

\bibitem[Arin and I\~narra(2001)]{ArinInarra2001}
Arin J, I\~narra E (2001) Egalitarian solutions in the core.
International Journal of Game Theory 30:187--193.

\bibitem[Arin et al.(2003)]{ArinKuipersVermeulen2003}
Arin J, Kuipers J, Vermeulen D (2003) Some characterizations of egalitarian
solutions on classes of TU-games. Mathematical Social Sciences 46:327--345.
\url{https://doi.org/10.1016/S0165-4896(03)00048-6}

\bibitem[Biswas et al.(1999)]{Biswas1999}
Biswas AK, Parthasarathy T, Potters JAM, Voorneveld M (1999) Large cores and
exactness. Games and Economic Behavior 28:1--12.
\url{https://doi.org/10.1006/game.1998.0683}

\bibitem[Cs\'oka et al.(2011)]{CsokaHeringsKoczy2011}
Cs\'oka P, Herings PJJ, K\'oczy L\'A (2011) Balancedness conditions for exact
games. Mathematical Methods of Operations Research 74:41--52.

\bibitem[Dietzenbacher(2021)]{Dietzenbacher2021}
Dietzenbacher B (2021) Monotonicity and egalitarianism.
Games and Economic Behavior 127:194--205.
\url{https://doi.org/10.1016/j.geb.2021.03.006}

\bibitem[Flanagan(2026)]{Flanagan2026}
Flanagan M (2026) Balanced transport at the egalitarian core: a complete
boundary phase diagram. Preprint; companion paper.

\bibitem[Llerena and Mauri(2017)]{LlerenaMauri2017}
Llerena F, Mauri L (2017) On the existence of the Dutta--Ray's egalitarian
solution. Mathematical Social Sciences 89:92--99.
\url{https://doi.org/10.1016/j.mathsocsci.2017.07.002}

\bibitem[Studen\'y and Kratochv\'il(2022)]{StudenyKratochvil2022}
Studen\'y M, Kratochv\'il V (2022) Facets of the cone of exact games.
Mathematical Methods of Operations Research 95:35--80.
\url{https://doi.org/10.1007/s00186-022-00770-4}

\end{thebibliography}

\end{document}
